\section{Anexos}
\begin{spacing}{1.5}

\subsection{Anexo A: reproducibilidad de los resultados}\label{anexo:reproducibilidad}

Todos los resultados numéricos de este documento se obtienen con el código del repositorio \texttt{tesis-generador-ecg}, y cada uno queda registrado en un archivo de salida junto con el entorno de cálculo con que se produjo.
% PENDIENTE: confirmar si el repositorio del código será público y agregar su URL.

\subsubsection{Entorno}

Los resultados se obtuvieron con Python 3.12.4, NumPy 2.1.3 y SciPy 1.16.0 en Windows, sobre un procesador Intel Core i7-7700HQ de cuatro núcleos. El archivo \texttt{scripts/generador/requirements.txt} fija esas versiones de NumPy y SciPy, porque L-BFGS-B converge a puntos ligeramente distintos entre versiones (Subsección~\ref{sec:mp}). La validación con una red neuronal requiere además PyTorch, y la medición independiente de intervalos, NeuroKit2. Las bases MIT-BIH, PTB-XL y PTB-XL+ se descargan de PhysioNet y no forman parte del repositorio.

El código incluye diez pruebas automatizadas (\texttt{pytest scripts/generador/tests}) que comprueban el gradiente analítico contra diferencias finitas, la equivalencia de la síntesis de la señal continua con la suma completa de los latidos, y la verificación y la calibración del motor de reglas.

\subsubsection{Correspondencia entre resultados y procedimientos}

La Tabla~\ref{tab:reproducibilidad} indica, para cada resultado, el procedimiento que lo produce y el archivo en que queda registrado. Las rutas son relativas a la carpeta \texttt{scripts/}.

\begin{table}[H]
\centering
\small
\resizebox{\textwidth}{!}{\begin{tabular}{|p{5.2cm}|p{5.4cm}|p{4.6cm}|}
\hline
\textbf{Resultado} & \textbf{Procedimiento} & \textbf{Archivo de salida} \\ \hline
Ajuste de los 2.237 latidos (Tablas~\ref{tab:metricas} y~\ref{tab:iteraciones}) & \path{generador/main_generator.py} \path{--procesos} \texttt{4} & \path{generador/dataset_v2_local.json} \\ \hline
Métricas de fidelidad & \path{generador/validate_performance.py} & \path{generador/metricas_v2_local.json} \\ \hline
Ablación de las correcciones (Tabla~\ref{tab:ablacion}) & \path{main_generator.py} con \path{--limites} y \path{--sin-linea-base} & \path{generador/metricas_*.json} \\ \hline
Segunda gaussiana para la onda T & \path{generador/experimento_onda_t.py} & \path{generador/experimento_onda_t.json} \\ \hline
Calibración con PTB-XL+ (Tabla~\ref{tab:calibracion}) & \path{diseno2/calibrar_rangos.py} y \path{diseno2/calibrar_plantilla.py} & \path{diseno2/resultados/factores_plantilla.json} \\ \hline
Cohorte etiquetada (Tabla~\ref{tab:cohorte}) & \path{generador/reglas_fisiologicas.py} \path{--verificar} & \path{generador/salida_cohorte/cohorte.json} \\ \hline
Cohorte sin calibrar, como referencia & \path{reglas_fisiologicas.py} \path{--sin-calibrar} \path{--salida} \path{salida_cohorte_sin_calibrar} & \path{generador/salida_cohorte_sin_calibrar/} \\ \hline
Medición independiente de intervalos (Tabla~\ref{tab:medicion_indep}) & \path{oe3/medicion_independiente.py} & \path{oe3/resultados/medicion_independiente.json} \\ \hline
Sesgo del estimador y cohorte de 24 horas (Tabla~\ref{tab:sesgo}) & \path{generador/validar_hurst.py} & \path{generador/validacion_hurst.json} \\ \hline
Análisis espectral (Tabla~\ref{tab:espectral}) & \path{oe5/analisis_espectral.py} y \path{oe5/espectro_cohorte.py} & \path{oe5/resultados/*.json} \\ \hline
Datos del diseño de validación con red neuronal & \path{diseno2/preparar_datos.py} \texttt{reales} y \path{sinteticos} & \path{salidas/diseno2/*.npz} \\ \hline
Red neuronal (Tablas~\ref{tab:red} y~\ref{tab:inversa}) & \path{diseno2/entrenar.py} \path{--semillas} \texttt{0} \texttt{1} \texttt{2} \path{--ablacion} & \path{diseno2/resultados/diseno2_resultados.json} \\ \hline
Régimen de pocos datos (Tabla~\ref{tab:fracciones}) & \path{entrenar.py} \path{--solo-fracciones} \texttt{0.01} \texttt{0.02} \texttt{0.05} \path{--semillas} \texttt{0} \texttt{1} \texttt{2} & \path{diseno2/resultados/diseno2_fracciones.json} \\ \hline
Evaluación con ECGFounder (Anexo~B) & \path{oe6/cohorte_calibrada_v2.py}, \path{oe6/inferir_ecgfounder.py} y \path{oe6/analizar_diseno1.py} & \path{oe6/resultados/diseno1_calibrado_v2_resultados.json} \\ \hline
\end{tabular}}
\caption{Procedimiento y archivo de salida de cada resultado reportado.}
\label{tab:reproducibilidad}
\end{table}

Los procedimientos con muestreo aleatorio reciben la semilla como parámetro y la registran en su salida: semilla 0 para la cohorte de verificación, 1 para la cohorte de entrenamiento de la red, 7 para la cohorte de veinticuatro horas, y 0, 1 y 2 para el entrenamiento de las redes. El ajuste morfológico no tiene componente aleatoria.

\subsection{Anexo B: evaluación exploratoria con un modelo preentrenado}\label{anexo:ecgfounder}

Antes de la validación con una red propia (Subsección~\ref{sec:met_red}), las señales sintéticas se evaluaron con ECGFounder \cite{li2025ecgfounder}, un modelo preentrenado sobre más de diez millones de electrocardiogramas y 150 etiquetas diagnósticas, en su versión de una derivación. Se le presentaron tiras de diez segundos a 500~Hz de los 1.000 sujetos de la cohorte etiquetada y 30 tiras reales del registro 100, con el preprocesado del repositorio oficial del modelo, sin reentrenarlo.

La evaluación tiene dos limitaciones que impiden usarla como validación principal. El modelo de una derivación se entrenó con la derivación I, mientras que la morfología del generador proviene de MLII, cercana a la derivación II. Como consecuencia, las etiquetas absolutas no son interpretables: el propio registro real se clasificó como anormal, con una probabilidad de 0,999. Por eso se analizaron solo comparaciones relativas: el AUROC de cada subgrupo con un rasgo desviado contra los sujetos sanos, sobre la probabilidad de la etiqueta diagnóstica correspondiente.

\begin{table}[H]
\centering
\resizebox{\textwidth}{!}{\begin{tabular}{|l|l|c|c|}
\hline
\textbf{Desvío} & \textbf{Etiqueta de ECGFounder} & \textbf{Cohorte sin calibrar} & \textbf{Cohorte calibrada} \\ \hline
PR desviado hacia abajo & PR corto & 0,983 & 0,988 \\ \hline
QRS desviado hacia arriba & QRS ancho & 0,751 & 0,891 \\ \hline
QT desviado hacia arriba & QT prolongado & 0,599 & 0,721 \\ \hline
PR desviado hacia arriba & bloqueo AV de primer grado & 0,505 & 0,614 \\ \hline
Onda P desviada hacia arriba & agrandamiento auricular izquierdo & 0,183 & 0,119 \\ \hline
QRS desviado hacia abajo (control) & QRS ancho & 0,017 & 0,028 \\ \hline
\end{tabular}}
\caption{AUROC de cada subgrupo contra los sujetos sanos en la probabilidad de la etiqueta de ECGFounder correspondiente, en la cohorte original y en la calibrada con PTB-XL+.}
\label{tab:ecgfounder}
\end{table}

Tres observaciones de esta evaluación orientaron el trabajo posterior. La primera es la que motivó la calibración de la Subsección~\ref{sec:calibracion}: en la cohorte sin calibrar, el PR desviado hacia arriba equivalía a unos 188~ms clínicos y la onda P ancha a unos 111~ms, valores normales, de modo que el modelo acertaba al no marcarlos; con la calibración, el AUROC de QRS ancho y de QT prolongado sube. La segunda es que el modelo sigue la duración del QRS en ambos sentidos ---el control con QRS angosto queda, como corresponde, en 0,028--- y el PR corto. La tercera es que, aun calibrados, los desvíos de PR largo y de onda P ancha no se reconocen: en una curva de dosis y respuesta con PR de 160 a 360~ms, la probabilidad de bloqueo AV de primer grado no se aparta de cero, y el modelo interpreta la onda P alejada del complejo como una alteración del ritmo. Esa observación anticipa el resultado de la validación con una red propia: el motor reproduce los intervalos, pero no los cambios de forma que en un registro real acompañan a esas alteraciones.

\end{spacing}

\newpage
